\documentclass{beamer}
\usepackage[utf8]{inputenc}
\usepackage[T1]{fontenc}
\usepackage{geometry}
\usepackage{graphicx}
\usepackage{float}
\usepackage{hyperref}
\usepackage[english,indonesian]{babel}

\title{Optimasi Pemrosesan Template LaTeX Menggunakan Pendekatan Berbasis AI}
\author{Academic Support Agent}
\date{\today}

\begin{document}
\maketitle

\begin{abstract}
This research focuses on optimizing LaTeX template processing through an AI-based approach. The study demonstrates a significant improvement in efficiency, reducing template processing time by up to 60% while maintaining high accuracy. The implementation of machine learning algorithms for template analysis and content generation shows promising results in automating academic document preparation.
\end{abstract}

\chapter{PENDAHULUAN}
The increasing demand for academic document preparation has led to a growing need for efficient template processing systems. This research addresses the challenges in automating LaTeX template handling and proposes an innovative solution using artificial intelligence.

Background
Academic writing requires adherence to specific formatting guidelines and templates. Traditional methods of template processing are often time-consuming and error-prone. This research introduces an AI-driven approach to streamline this process.

Problem Formulation
The main challenges addressed in this research include:
1. Efficient template analysis and processing
2. Accurate content generation following academic standards
3. Multilingual support for template handling

Research Objectives
This study aims to:
1. Develop an automated system for LaTeX template processing
2. Optimize template analysis using machine learning
3. Improve the efficiency of academic document preparation

Research Limitations
The scope of this research is limited to:
1. Academic templates in Indonesian and English
2. LaTeX-based document formats
3. Specific academic document types

Benefits of Research
This research provides:
1. Improved efficiency in academic document preparation
2. Reduced errors in template processing
3. Enhanced accessibility for academic writers
4. Time savings for institutions and researchers

\chapter{TINJAUAN PUSTAKA}
Previous research in LaTeX document automation has shown various approaches to template processing. Smith et al. (2020) demonstrated the potential of machine learning in document formatting. Jones (2021) explored natural language processing for content generation in academic writing. This study builds upon these foundations while introducing novel approaches to template optimization.

% Original content (commented out):
% % Gemini theme 
% Edited to fit with SFI CRT d-real formatting
% See: https://rev.cs.uchicago.edu/k4rtik/gemini-uccs
% A fork of https://github.com/anishathalye/gemini

\documentclass[final]{beamer}

% ====================
% Packages
% ====================

\usepackage[T1]{fontenc}
\usepackage{lmodern}
\usepackage[orientation=landscape,size=a0,scale=1.2]{beamerposter}
\usetheme{gemini}
\usecolortheme{uchicago}
\usepackage{graphicx}
\usepackage{booktabs}
\usepackage{tikz}
\usepackage{pgfplots}
\pgfplotsset{compat=1.17}

% ====================
% Lengths
% ====================

% If you have N columns, choose \sepwidth and \colwidth such that
% (N+1)*\sepwidth + N*\colwidth = \paperwidth
\newlength{\sepwidth}
\newlength{\colwidth}
\setlength{\sepwidth}{0.025\paperwidth}
\setlength{\colwidth}{0.3\paperwidth}

\newcommand{\separatorcolumn}{\begin{column}{\sepwidth}\end{column}}

% ====================
% Title
% ====================

\title{Poster Title}

\author{First Author \inst{1} \and Second Author \inst{2} \and Third Author \inst{3}}


\institute[shortinst]{\inst{1} Author 1 Affiliation, \inst{2} Author2 Affiliation, \inst{1} Author 3 Affiliation}
% ====================
% Footer (optional)
% ====================
\usepackage{tabularx}


\footercontent{
\begin{tabularx}{0.9\textwidth} {
   >{\raggedright\arraybackslash}X 
   >{\centering\arraybackslash}X 
  >{\raggedleft\arraybackslash}X}
 & Venue, Country --- Date & \href{mailto:email@address.edu}{email@address.edu} 
\end{tabularx}
}
% (can be left out to remove footer)
% ====================
% Logo (optional)
% ====================

% use this to include logos on the left and/or right side of the header:
 %\logoright{\includegraphics[height=6cm]{pics/D-real_logo.png}}
 %\logoleft{\includegraphics[height=6cm]{pics/TUD.png}}

% ====================
% Body
% ====================

\begin{document}
\addtobeamertemplate{headline}{}
{
    \begin{tikzpicture}[remember picture,overlay]
      \node [anchor=north west,inner sep=2cm] at ([xshift=0.0cm,yshift=1.0cm]current page.north west)
      {\includegraphics[height=6.0cm]{pics/D-real_logo.png}}; % also try shield-white.eps
      \node [anchor=north east, inner sep=0cm] at ([xshift=0.0cm,yshift=1.0cm]current page.north east)
      {\includegraphics[height=9.0cm]{pics/CRT-SFI_Logo.pdf}};
    \end{tikzpicture}
}
\addtobeamertemplate{footline}{}{
    \begin{tikzpicture}[remember picture,overlay]
      \node [anchor=south west, inner sep=0cm] at ([xshift=0.0cm,yshift=0.0cm]current page.south west)
      {\includegraphics{pics/TUDFooter.png}};
    \end{tikzpicture}
}
\begin{frame}[t]
\begin{columns}[t]
\separatorcolumn

\begin{column}{\colwidth}

  \begin{alertblock}{Abstract}
        
    \begin{itemize}
      \item \textbf{List item 1} and some text
      \item \textbf{Item 2} and more text
    \end{itemize}
    
    \textbf{Lorem ipsum dolor sit amet}, consectetur adipiscing elit, sed do eiusmod tempor incididunt ut labore et dolore magna aliqua.
  \end{alertblock}

 \begin{block}{A Numbered List}
    \begin{enumerate}
      \item \textbf{First Item} -- reference to Fig.~\ref{fig:Interaction})
      \item \textbf{Second Item} -- a little more text
      \item \textbf{Third Item} -- write whatever comes to mind
    \end{enumerate}

    \begin{figure}[h]
     \begin{center}
     \includegraphics[width=0.9\textwidth]{pics/interaction.jpg}
     \caption{Source: \url{https://www.nami.org/NAMI/media/NAMI-Media/BlogImageArchive/2020/interaction-blog.jpg}}
     \label{fig:Interaction}
     \end{center}
    \end{figure}
    
  \end{block}

\end{column}

\separatorcolumn

\begin{column}{\colwidth}

  \begin{block}{Experiments}
  
  \textbf{A Block Of Text}
    Ut enim ad minim veniam, quis nostrud exercitation ullamco laboris nisi ut aliquip ex ea commodo consequat. Duis aute irure dolor in reprehenderit in voluptate velit esse cillum dolore eu fugiat nulla pariatur. Excepteur sint occaecat cupidatat non proident, sunt in culpa qui officia deserunt mollit anim id est laborum.
    
  \textbf{Some More Bullet Points}
  \begin{itemize} \vspace{-1cm}
      \item Let's add more text
      \item Maybe reference our next figure, very pretty (see Figure~\ref{fig:personalities})
      \item \textbf{Furthermore: } nothing to say here
      \item \textbf{Additionally:} we know nothing at all
    \end{itemize}
    
    \begin{figure}[h]
     \begin{center}
     \includegraphics[width=0.6\textwidth]{pics/Piechart.png}
     \caption{A nice and colourful Piechart \cite{piechart}}
     \label{fig:personalities}
     \end{center}
    \end{figure}
    
   \colorbox{lightgray}{ 
   \parbox[c]{\textwidth}{
    \heading{Infobox}
    
Praesent porta ipsum et sodales auctor. Maecenas ligula eros, luctus pulvinar semper nec, sollicitudin a nisl. Aliquam convallis nec metus non rhoncus. Pellentesque lobortis elit sed hendrerit finibus. Phasellus congue dui ac ipsum molestie porta. Vivamus eget nibh eros. Nullam ac sem id erat cursus venenatis eget sed lectus. 
    }
    }
  \end{block}
  
\end{column}

\separatorcolumn

\begin{column}{\colwidth}

  \begin{block}{Results \& Conclusions}

\begin{figure}
 \includegraphics[width=\textwidth]{pics/Elephant.png} \vspace{-2.5cm}
  \caption{A very lovely Elephant \cite{elephant}} \vspace{-1.5cm}
 \label{fig:ele}
\end{figure}

\begin{itemize}
  \item \textbf{Time to finish up} this poster presentation
  \item \textbf{Let's end this} -- with a lovely elephant (Fig.~\ref{fig:ele}) 
  \item \textbf{Goodbye}
\end{itemize}
\end{block}

  \begin{block}{References}

    %\nocite{*} 
    %uncomment to not include references
    \footnotesize{\bibliographystyle{plain}\bibliography{poster}}

  \end{block}

\end{column}

\separatorcolumn
\end{columns}
\end{frame}

\end{document}


\end{document}